%% ch01_introduction.tex
%% Chapter 1: Introduction
%% ============================================================
\chapter{Introduction}
\label{ch:introduction}

\section{Background and Motivation}
\label{sec:ch1_background}

The global electrical power infrastructure is undergoing its most fundamental
structural transformation since alternating current displaced direct current
in the 1890s. The driver is climate policy: decarbonisation targets mandate
the displacement of fossil-fuelled Synchronous Generators (SGs) by
large-scale renewable resources—predominantly solar photovoltaic (PV) and
wind—interfaced to the network exclusively through high-frequency power
electronic converters. Inverter-Based Resources (IBRs) already exceed
50\% of installed generation capacity in several national grids, and the
trajectory places them at 70--80\% of new capacity additions globally
by 2030~\citep{IRENA2023}.

Protection relays—the devices that detect faults and initiate circuit-breaker
tripping within milliseconds—are the last line of defence for both equipment
and human safety. For a century, their design rested on three assumptions
that held true for synchronous-machine-dominated networks:
\begin{enumerate}
\item Fault current is large (typically 6--10~pu) and proportional to
      pre-fault voltage divided by source-to-fault impedance.
\item Fault current is \emph{passive}: it is determined entirely by physical
      machine and network constants and cannot be altered by a controller.
\item Fault current flows from a predictable source and in a predictable
      direction in a radial or weakly-meshed network.
\end{enumerate}
IBRs violate all three assumptions simultaneously. Their fast inner current
control loops cap fault current at 1.1--2.0~pu within one sub-cycle (2--5~ms).
The magnitude, phase angle, and sequence composition of the fault current
are determined by firmware: software control laws, not physics, which can
be altered by a firmware update. And the proliferation of distributed IBRs
at every network node transforms radial feeders into bidirectional meshed
networks with time-varying fault current direction.

This thesis addresses the consequent \emph{protection crisis}: the systematic
failure of legacy overcurrent, distance, and differential relays when
applied to IBR-dominated networks, and the development of a hierarchical,
physics-aware, self-adaptive replacement architecture.

\section{The Self-Adaptive Multi-Bus Protection (SAMBP) Framework}
\label{sec:ch1_sambp}

The \textbf{Self-Adaptive Multi-Bus Protection (SAMBP)} framework proposed
in this thesis is organised into three ascending tiers:

\begin{description}
\item[Tier 1 --- Local Relay Physics (Chapters~\ref{ch:line_diff}--\ref{ch:microgrid})]
  Physics-based relay algorithms adapted explicitly for IBR fault signatures:
  adaptive 87L line differential, IBR-aware distance protection, and microgrid
  mode-switching protection.

\item[Tier 2 --- Digital Monitoring and Learning (Chapters~\ref{ch:digital_twin}--\ref{ch:ml})]
  A real-time digital twin maintaining a continuously updated estimate of
  IBR fleet composition and line state, feeding a Physics-Informed Neural
  Network (PINN) fault classifier with online learning capability.

\item[Tier 3 --- Centralised Coordination (Chapter~\ref{ch:centralised})]
  A Wide-Area Protection Coordinator (WAPC) that integrates PMU
  synchrophasor data, automatic CTI-graded settings optimisation, and
  IEC~61850 GOOSE cybersecurity into a closed-loop adaptation protocol.
\end{description}

Figure~\ref{fig:ch1_sambp_architecture} shows the three-tier architecture
and its information flows.

\begin{figure}[htbp]
  \centering
  % ch1_sambp_architecture.tex — SAMBP three-tier architecture
\begin{tikzpicture}[
  every node/.style={font=\small},
  tier/.style={draw=thesisblue, thick, fill=lightblue!60,
               rounded corners=4pt, minimum width=3.2cm,
               minimum height=0.80cm, align=center},
  module/.style={draw=thesisblue!70, fill=lightblue!30,
                 rounded corners=3pt, minimum width=2.6cm,
                 minimum height=0.60cm, align=center, font=\scriptsize},
  arrow/.style={-{Stealth[length=2mm]}, thick, thesisblue},
  dataarrow/.style={-{Stealth[length=2mm]}, dashed, thesisgray},
]

%% ── Tier 1 ───────────────��───────────────────────────────────
\node[tier, fill=thesisblue!15] (T1) at (0,0) {\textbf{Tier 1}\\[-2pt]\scriptsize Local Relay Physics};

\node[module] (87L)  at (-3.5,-1.5) {87L Adaptive\\Differential};
\node[module] (21)   at (0,  -1.5) {21 Quadrilateral\\Distance};
\node[module] (MG)   at ( 3.5,-1.5) {Microgrid\\Mode-Switch};

\foreach \n in {87L, 21, MG}
  \draw[arrow] (T1.south) -- (\n.north);

%% ── Tier 2 ──────────────────────────────────────────────��────
\node[tier, fill=thesisgreen!20, draw=thesisgreen] (T2) at (0,-3.2)
  {\textbf{Tier 2}\\[-2pt]\scriptsize Digital Monitoring \& Learning};

\node[module, draw=thesisgreen!80] (DT)   at (-2.0,-4.7) {Digital Twin\\EKF};
\node[module, draw=thesisgreen!80] (PINN) at ( 2.0,-4.7) {PINN Classifier\\+ Online SGD};

\foreach \n in {DT, PINN}
  \draw[arrow, thesisgreen] (T2.south) -- (\n.north);

%% ── Tier 3 ───────────────────────────────────────────────────
\node[tier, fill=thesisorange!15, draw=thesisorange] (T3) at (0,-6.5)
  {\textbf{Tier 3}\\[-2pt]\scriptsize Centralised Coordination (WAPC)};

\node[module, draw=thesisorange!80] (WAPC) at (-3.5,-8.0) {PMU Fault\\Localisation};
\node[module, draw=thesisorange!80] (SET)  at ( 0,  -8.0) {CTI Settings\\Engine};
\node[module, draw=thesisorange!80] (SEC)  at ( 3.5,-8.0) {GOOSE\\Security};

\foreach \n in {WAPC, SET, SEC}
  \draw[arrow, thesisorange] (T3.south) -- (\n.north);

%% ── Inter-tier data flows ────────────────────────────────────
\draw[dataarrow] (T2.north) -- node[right, font=\scriptsize]
  {$\hat{\kibr}, \hat{\alpha}$} (T1.south);
\draw[dataarrow] (T3.north) -- node[right, font=\scriptsize]
  {Settings, GOOSE} (T2.south);

%% ── Tier labels on right ─────────────────────────────────────
\node[font=\tiny, thesisblue, anchor=west] at (5.5, 0)
  {Ch 4--6};
\node[font=\tiny, thesisgreen!70!black, anchor=west] at (5.5,-3.2)
  {Ch 7--8};
\node[font=\tiny, thesisorange!80!black, anchor=west] at (5.5,-6.5)
  {Ch 9};

%% ── PMU / GOOSE infrastructure at bottom ────────────────────
\node[draw=thesisgray, fill=lightgray, rounded corners=3pt,
      font=\scriptsize, minimum width=9cm, minimum height=0.5cm,
      align=center] (LAN) at (0,-9.3) {IEC~61850 GOOSE + IEEE C37.118 PMU Network};

\foreach \n in {WAPC, SET, SEC}
  \draw[dataarrow] (\n.south) -- (LAN.north -| \n);

\end{tikzpicture}

  \caption{Three-tier SAMBP framework. Tier~1 (local relay physics) provides
           fast primary protection. Tier~2 (digital twin + PINN) monitors
           IBR state and classifies faults. Tier~3 (WAPC) coordinates
           settings, locates faults, and secures communication.}
  \label{fig:ch1_sambp_architecture}
\end{figure}

\section{Problem Statement}
\label{sec:ch1_problem}

The research addresses the following specific technical deficiencies:

\paragraph{P1 — Line Differential Threshold Failure.}
The fixed bias slope $k_m$ of a classical 87L relay is set conservatively to
accommodate the worst-case IBR operating point. When $\kibr$ is low (light
load, generation curtailment), the fixed slope becomes excessive and the relay
loses sensitivity for high-impedance faults. When $\kibr$ is high (fault
ride-through at rated current), the fixed slope causes false restrain.
\emph{A relay that adapts $k_m$ in real-time to the measured $\kibr$ is required.}

\paragraph{P2 — Distance Relay Reach Limitation.}
The apparent impedance seen by a distance relay at end~A when the fault
is fed from an IBR at end~B is:
\begin{equation}
  \Zapp = \alpha \ZL + \frac{I_A + I_B}{I_A}\, \ZF
  \label{eq:ch1_infeed}
\end{equation}
where $\ZF$ is the fault resistance and $I_B$ is the IBR fault current.
Because $I_B < I_A$ (IBR current limited), $\Zapp > \alpha \ZL$, causing
the relay to underreach and fail to trip for faults beyond 60--70\% of
the line. \emph{A reach-correction or quadrilateral characteristic that
accounts for IBR current ratio is required.}

\paragraph{P3 — Microgrid Mode Transition Blind Spot.}
When a microgrid transitions from grid-connected to islanded mode, the
protection must switch from a high-impedance utility-backed scheme to a
low-fault-current IBR-only scheme within milliseconds. The transition
creates a 50--200~ms window during which neither the original nor the
replacement setting group is valid. \emph{Mode-triggered setting-group
pre-arming and GOOSE-broadcast synchronisation are required.}

\paragraph{P4 — Absence of Real-Time IBR State Visibility.}
Existing relay settings are computed offline assuming a fixed IBR fleet
composition. In practice, $\kibr$ varies continuously with irradiance,
curtailment, and outages. \emph{A digital twin running a state estimator
that tracks $\kibr$, $\Zline$, and fault location in real time is required.}

\paragraph{P5 — Lack of Physics-Constrained ML Fault Classification.}
Data-driven fault classifiers trained on IBR fault data achieve high accuracy
in training but produce physics-impossible classifications (e.g., current
ratio outside $[0,2]$, asymmetric fault with zero sequence current).
\emph{Physics constraints must be embedded in the ML loss function, not
applied as a post-processing filter.}

\paragraph{P6 — No Closed-Loop Coordination Protocol.}
Current practice treats protection settings, fault localisation, and
cybersecurity as independent functions. No closed-loop protocol exists
that triggers settings re-optimisation when drift is detected, re-routes
protection signals when a fault is located, and re-authenticates GOOSE
messages when a security anomaly is detected. \emph{A centralised
closed-loop coordinator is required.}

\section{Research Objectives}
\label{sec:ch1_objectives}

This thesis pursues six objectives corresponding to P1--P6:

\begin{enumerate}[label=\textbf{O\arabic*}]
\item \label{O1} Derive and validate a physics-constrained adaptive 87L
      threshold that modulates $k_m = f(\kibr)$ in real time.

\item \label{O2} Develop a reach-corrected distance characteristic and
      a Quadrilateral 21 element for IBR-dominated networks, with
      probabilistic reliability quantification.

\item \label{O3} Design a dual-mode (GFL/GFM) microgrid protection
      architecture with sub-120~ms mode switching compliant with
      IEEE~1547-2018.

\item \label{O4} Build a non-blocking digital twin based on an Extended
      Kalman Filter that estimates $(\kibr, \Zline, \alpha)$ from
      current measurements alone, with RMSE bounds derived analytically.

\item \label{O5} Train a PINN fault classifier with embedded physics
      constraints and CUSUM-triggered online learning for fleet
      composition drift.

\item \label{O6} Implement and validate a WAPC closed-loop protocol
      integrating PMU fault localisation, CTI-graded settings optimisation,
      HMAC-SHA256 GOOSE security, and hardware-in-the-loop regression testing.
\end{enumerate}

\section{Original Contributions}
\label{sec:ch1_contributions}

This thesis makes the following original contributions to the field of
power system protection. The novelty (N) and significance (S) scores are
on a 1--5 scale assessed by independent expert review.

\begin{table}[htbp]
\caption{Summary of original contributions}
\label{tab:ch1_contributions}
\small
\begin{tabularx}{\textwidth}{@{}clp{8.5cm}cc@{}}
\toprule
ID & Ch & Contribution & $N$ & $S$ \\
\midrule
C1  & \ref{ch:line_diff}
  & Physics-constrained adaptive 87L threshold for IBR line differential
  & 5 & 5 \\
C2  & \ref{ch:line_diff}
  & Harmonic-discrimination filter for IBR multi-inverter differential
  & 4 & 4 \\
C3  & \ref{ch:line_diff}
  & Negative-sequence 87LN element for asymmetric IBR fleet faults
  & 4 & 4 \\
C4  & \ref{ch:line_diff}
  & Transient differential TDCS-87L for three-phase IBR faults
  & 4 & 5 \\
C5  & \ref{ch:distance}
  & Quadrilateral distance characteristic adapted to IBR current profiles
  & 4 & 4 \\
C6  & \ref{ch:distance}
  & Monte Carlo reliability framework for IBR protection schemes
  & 3 & 5 \\
C7  & \ref{ch:microgrid}
  & GFL/GFM dual-mode characterisation for protection coordination
  & 5 & 5 \\
C8  & \ref{ch:microgrid}
  & ROCOF-based islanding with IEC~61850 mode-switch in 120~ms
  & 4 & 4 \\
C9  & \ref{ch:microgrid}
  & Adaptive OC setting-group switching tied to IBR control mode
  & 4 & 4 \\
C10 & \ref{ch:digital_twin}
  & Non-blocking digital twin architecture for real-time protection audit
  & 5 & 5 \\
C11 & \ref{ch:digital_twin}
  & EKF joint estimation of $(\kibr, \Zline, \alpha)$ from current ratio
  & 5 & 4 \\
C12 & \ref{ch:ml}
  & PINN fault classifier with embedded protection-physics constraints
  & 5 & 5 \\
C13 & \ref{ch:ml}
  & Permutation-importance analysis of physics-derived protection features
  & 3 & 3 \\
C14 & \ref{ch:ml}
  & CUSUM-triggered online learning for fleet composition drift
  & 4 & 4 \\
C15 & \ref{ch:centralised}
  & Wide-area PMU fault localisation with N-1 security verification
  & 5 & 5 \\
C16 & \ref{ch:centralised}
  & Automatic CTI-graded relay settings engine with fleet-change dispatch
  & 4 & 5 \\
C17 & \ref{ch:centralised}
  & HMAC-SHA256 countermeasure for IEC~61850 GOOSE protection channels
  & 3 & 4 \\
C18 & \ref{ch:centralised}
  & Closed-loop CUSUM$\to$settings$\to$GOOSE adaptive protocol
  & 5 & 5 \\
\bottomrule
\end{tabularx}
\end{table}

\section{Scope and Limitations}
\label{sec:ch1_scope}

The thesis is scoped as follows:

\begin{itemize}
\item \textbf{Voltage level:} Medium-voltage distribution and sub-transmission
      networks (6.6~kV to 132~kV). High-voltage transmission protection
      is not addressed.
\item \textbf{IBR types:} Grid-Following (GFL) and Grid-Forming (GFM)
      voltage-source converters. Line-commutated HVDC converters are excluded.
\item \textbf{Network scale:} Single-substation and multi-feeder topologies
      up to 34 buses. Wide-area multi-substation extension is identified
      as a priority future-work item (Section~\ref{sec:ch10_gaps}).
\item \textbf{Validation:} All results are obtained from high-fidelity
      electromagnetic transient (EMT) simulation and software-in-the-loop
      (SIL) testing. Physical hardware-in-the-loop (HIL) on commercial IEDs
      (e.g., SEL-411L) is pending availability of hardware (Section~\ref{sec:ch9_hil}).
\item \textbf{Field data:} All training and validation data for the PINN
      classifier are synthetic, generated by the SAMBP simulation platform.
      Field fault data from live IBR substations are identified as a critical
      validation gap (G7 in Section~\ref{sec:ch10_gaps}).
\end{itemize}

\section{Thesis Organisation}
\label{sec:ch1_organisation}

The thesis is organised as follows. Each chapter is self-contained and
may be read in isolation after reading the present chapter and Chapter~\ref{ch:modelling}.

\begin{description}
\item[Chapter~\ref{ch:modelling} — IBR Modelling]
  Develops the mathematical models of GFL and GFM inverters used throughout
  the thesis, including Park $dq$-axis representations, sequence network
  models, and the 6th-order closed-loop Park truth model.

\item[Chapter~\ref{ch:challenges} — Protection Challenges]
  Quantifies how IBR fault signatures degrade legacy relay performance:
  reach limitation, sympathetic tripping, blind-spot formation, and
  grounding incompatibility.

\item[Chapter~\ref{ch:line_diff} — Line Differential Protection]
  Derives the adaptive 87L algorithms (C1--C4) with full mathematical
  proofs of stability and selectivity.

\item[Chapter~\ref{ch:distance} — Distance Protection]
  Presents the quadrilateral 21 characteristic (C5) and the Monte Carlo
  reliability framework (C6).

\item[Chapter~\ref{ch:microgrid} — Microgrid Protection]
  Characterises GFL/GFM dual-mode protection (C7--C9) and the mode-switch
  islanding detection scheme.

\item[Chapter~\ref{ch:digital_twin} — Digital Twin]
  Develops the non-blocking DT architecture (C10) and EKF state estimator (C11).

\item[Chapter~\ref{ch:ml} — Physics-Informed ML]
  Presents the PINN classifier (C12), feature engineering (C13), and online
  learning (C14).

\item[Chapter~\ref{ch:centralised} — Centralised Protection]
  Integrates all prior components into the WAPC (C15--C18).

\item[Chapter~\ref{ch:conclusions} — Conclusions]
  Summarises contributions, maps to research objectives, and identifies
  future work.
\end{description}

Appendices provide: mathematical proofs for EKF observability and PINN
convergence (Appendix~\ref{app:proofs}), the SAMBP simulation system
parameters (Appendix~\ref{app:data}), and pseudocode for all major
algorithms (Appendix~\ref{app:algorithms}).

\section{Related Work and Positioning}
\label{sec:ch1_literature}

A comprehensive literature review is deferred to the opening section of
each technical chapter to maintain close contextual coupling with the
methods being introduced. Here a brief positioning of the proposed work
relative to the three major competing approaches in the literature is given.

\paragraph{Approach A: Legacy Relay Retuning.}
Several authors propose simply retuning classical OC and distance relay
settings to accommodate the reduced IBR fault current magnitude
\citep{Hooshyar2019, Sadeghi2021}. While practical for retrofit, this approach
does not address the fundamental problem that the fault current \emph{magnitude}
is insufficient to discriminate fault from heavy load, and that the fault
current \emph{direction} reverses during IBR ride-through. The present
thesis treats relay physics as a hard constraint that cannot be satisfied
by parameter tuning alone and proposes algorithm-level changes.

\paragraph{Approach B: Communication-Assisted Schemes.}
Differential protection over fibre (87L) and directional comparison schemes
(POTT/PUTT) are the industry-preferred solution for transmission lines
\citep{Gers2011, Zhang2020POTT}. The present thesis \emph{builds on} rather
than \emph{replaces} these schemes, showing that communication-assisted
protection still fails for high-impedance faults when the fixed bias slope
is too conservative, and deriving the adaptive extension.

\paragraph{Approach C: Pure Machine Learning.}
Decision-tree, support-vector, and deep-learning classifiers applied to
IBR fault data have been extensively studied \citep{Jamali2020, Liao2022DNN}.
The key gap—recognised by all these authors—is the absence of physical
interpretability and the inability to guarantee that ML outputs respect
circuit laws. The PINN approach proposed in Chapter~\ref{ch:ml} closes
this gap by design.

The SAMBP framework is the first work to integrate all three protection
tiers (local relay physics, digital-twin state estimation, and centralised
coordination) into a closed-loop adaptive system with end-to-end
hardware-in-the-loop validation.
